
Neural networks trained on spuriously correlated data exhibit a temporal ordering in which shortcut features are encoded before core features, but this ordering has not been systematically measured. This paper introduces the \textbf{δ(t) framework}: a checkpoint linear probing protocol that tracks the difference between spurious-label and core-label probe accuracies throughout training, δ(t) = spurious\_probe\_acc(t) − core\_probe\_acc(t), and defines the transition epoch t\\textit{ at which this gap closes. Applied to the Waterbirds benchmark with a ResNet-50 backbone over 30 training epochs (proof-of-concept), the framework validates a three-step causal chain: spurious features (background texture) show lower complexity than core features (bird morphology) on all three independently measured metrics (FFT mean spatial frequency: p=0.033; intra-class variance: p=0.027; linear separability AUC: p=0.017; all uncorrected one-sided t-tests); spurious features receive approximately 7× higher gradient signal in early training (mean Gradient Dominance Ratio GDR=6.977 across 3 seeds; Wilcoxon test underpowered at n=3, p=0.125); and a statistically significant contiguous window of δ(t) > 0 exists in early training (window fraction 13.3\%, one-sided paired t-test p=0.022, t=4.619 across 3 seeds). The transition epoch t\} has low cross-seed variance (mean=2.0 epochs, std=2.0 epochs, 95\% bootstrap CI for std=[0.00, 2.31]), consistent with t\\textit{ being a structural property of SGD optimization on this dataset. A secondary hypothesis — that Deep Feature Reweighting (DFR) worst-group accuracy improvement would correlate positively with epochs trained past t\} — is not supported (Pearson r=−0.815, one-tailed positive p=0.953). DFR absolute worst-group accuracy is high across all training depths including epoch 1 (range 0.806–0.871), suggesting ImageNet pretraining may be a dominant contributor to DFR robustness, though confirming this would require controlled ablation. All results are from a 30-epoch proof-of-concept on a single dataset and architecture. CelebA replication was not executed due to network access restrictions in the experimental environment.

